\documentclass[10pt,a4paper]{amsart}
\usepackage[margin=19mm]{geometry}
\usepackage{amsmath,amssymb,mathtools}
\usepackage[scr=rsfs,cal=euler]{mathalfa}
\usepackage{xfrac}
\usepackage[unicode,hidelinks,bookmarks=false]{hyperref}
\newcommand{\ie}{i.e.\ }
\newcommand{\cf}{cf.\ }
\newcommand{\resp}{respectively\ }
\newcommand{\Subsection}{\subsection}
\providecommand{\nobreakdash}{\nobreak\hbox{-}}
\setlength{\emergencystretch}{2em}
\multlinegap0pt
\usepackage{bm}
\usepackage{bbm}
\usepackage{dsfont}
\usepackage{xspace}
\usepackage{cancel}
\usepackage{mathtools}
\usepackage{amsthm}
\makeatletter
\@ifundefined{theorem}{\newtheorem{theorem}{Theorem}}{}
\@ifundefined{definition}{\newtheorem{definition}[theorem]{Definition}}{}
\@ifundefined{proposition}{\newtheorem{proposition}[theorem]{Proposition}}{}
\makeatother
\title[Spencer Differential Degeneration Theory and Its Application]{Spencer Differential Degeneration Theory and Its Applications in Algebraic Geometry}
\author{Dongzhe Zheng}
\date{Source: arXiv:2506.07410v2 (CC BY 4.0)}
\begin{document}
\maketitle

\noindent\emph{Source and licence.} Spencer Differential Degeneration Theory and Its Applications in Algebraic Geometry. Version arXiv:2506.07410v2 (CC BY 4.0), distributed under the Creative Commons Attribution 4.0 International licence, as recorded in the arXiv licence metadata for this exact version and shown on the abstract page https://arxiv.org/abs/2506.07410v2. Full rights notice: licenses/arxiv-2506.07410-CC-BY-4.0.txt.\par\smallskip

\noindent\textit{Editorial scope.} Counted unit: the Theorem whose source label is \texttt{thm:composite\_map\_structure\_and\_surjectivity} in the article (source0.tex lines 626--634, source label \texttt{thm:composite\_map\_structure\_and\_surjectivity}), delivered together with the article's own complete original proof of it (source0.tex lines 636--651, one \texttt{proof} environment of 16 lines). No mathematics is written, repaired or re-derived: statement and proof are the article's own text line for line. Required context retained uncounted: prop:projection\_basic\_properties (lines 608--611); def:composite\_map\_to\_11 (lines 619--624). Every definition, display and label of those lines is retained unchanged. Omitted from this package and not redistributed: 1--607, 612--618, 625--625, 635--635, 652--799. Nothing inside a retained excerpt is omitted or altered: every retained body is delivered exactly as the article's lines are written, so no omission receipt of any kind accompanies this package. Citations: the retained excerpts carry no citation command at all, so no bibliography entry of the article is transcribed and no reference is redistributed. Reference adaptation: none was needed and none was made; every reference inside the delivered unit resolves inside it, so no unresolved reference or citation is produced. Printed slips kept literal: no printed word, symbol, hypothesis or number of the counted statement or of its proof was altered. Layout and class: the article's own class is \texttt{extarticle} with packages including inputenc, geometry, graphicx, booktabs, mathrsfs, bm, xcolor, algorithm, algpseudocode, mathtools, enumitem, tikz-cd, newtxtext, newtxmath; the wrapper is the lane's pinned \texttt{amsart} editorial wrapper at 10pt on A4, which restates the theorem-like environments the retained text needs and adds the article's own macro definitions that the retained text needs (none), with extra wrapper packages (bm, bbm, dsfont, xspace, cancel, mathtools). The delivered numbering is the unit's own. Assets: no retained span contains a figure environment, a graphics-inclusion command, a PDF-page inclusion, a table or a TikZ picture; no image file is copied into this package, the package declares no asset dependency and no image file is redistributed with it. Licence: version arXiv:2506.07410v2 of this article is distributed under the Creative Commons Attribution 4.0 International licence, as recorded in the arXiv licence metadata for this exact version and shown on the abstract page https://arxiv.org/abs/2506.07410v2. A full rights notice is delivered at licenses/arxiv-2506.07410-CC-BY-4.0.txt. Characters: every retained excerpt is pure ASCII, so the delivered engine, XeLaTeX with its default Latin Modern text font, reports no missing character.
\begin{proposition}[Basic Properties of Projection Mapping]
\label{prop:projection_basic_properties}
The mapping $\Pi^k$ is a well-defined linear surjection.
\end{proposition}

\begin{definition}[Composite Mapping to $(1,1)$-Cohomology]
\label{def:composite_map_to_11}
Define the composite mapping:
$$\Phi_{D,\lambda}: Z_{deg}^2(D,\lambda) \xrightarrow{\Pi^2} Z^2_{dR}(X) \hookrightarrow H^2_{dR}(X) \xrightarrow{\pi^{1,1}} H^{1,1}(X, \mathbb{C})$$
where $\pi^{1,1}$ is the standard Hodge projection.
\end{definition}

\begin{theorem}[Structure and Surjectivity of Composite Mapping]
\label{thm:composite_map_structure_and_surjectivity}
On K3 surface $X$, the composite mapping $\Phi_{D,\lambda}: Z_{deg}^2(D,\lambda) \to H^{1,1}(X, \mathbb{C})$ has the following structural properties:
\begin{enumerate}
    \item $\Phi_{D,\lambda}$ is a well-defined linear mapping.
    \item $\Phi_{D,\lambda}$ is a \textbf{surjection}, whose image space is the entire $H^{1,1}(X, \mathbb{C})$.
    \item Therefore, the dimension of its image space is constantly $\dim(\text{im}(\Phi_{D,\lambda})) = h^{1,1}(X) = 20$.
\end{enumerate}
\end{theorem}

\begin{proof}
\begin{enumerate}
    \item \textbf{Well-definedness}: The mapping is defined as a composition of well-defined linear mappings (as shown in Definition \ref{def:composite_map_to_11}), therefore it is itself a well-defined linear mapping.

    \item \textbf{Surjectivity}: We analyze each mapping in the composition chain:
    $$ Z_{deg}^2(D,\lambda) \xrightarrow{\Pi^2} Z^2_{dR}(X) \xrightarrow{q} H^2_{dR}(X, \mathbb{C}) \xrightarrow{\pi^{1,1}} H^{1,1}(X, \mathbb{C}) $$
    \begin{itemize}
        \item The mapping $\Pi^2: Z_{deg}^2(D,\lambda) \to Z^2_{dR}(X)$ is surjective (proven in Proposition \ref{prop:projection_basic_properties}), provided that $\mathcal{K}_{\mathfrak{su}(2)}^2(\lambda)$ is nonzero.
        \item The quotient mapping $q: Z^2_{dR}(X) \to H^2_{dR}(X, \mathbb{C})$ is surjective by definition.
        \item For K3 surfaces, the Hodge projection $\pi^{1,1}: H^2_{dR}(X, \mathbb{C}) \to H^{1,1}(X, \mathbb{C})$ is also surjective, since both the source and target spaces are nonzero.
    \end{itemize}
    As a composition of surjections, $\Phi_{D,\lambda}$ is necessarily surjective.

    \item \textbf{Dimension conclusion}: Since the mapping is surjective, its image space is the entire target space $H^{1,1}(X, \mathbb{C})$. For projective K3 surfaces, we know their Hodge number $h^{1,1}(X)=20$. Therefore the conclusion holds.
\end{enumerate}
\end{proof}

\end{document}
